%====================================================================
%  WIP_intro.tex  --  introduction for the Sharp_NV paper.
%
%  Editorial passes (prose-divergence from KdV companion \cite{chcora}):
%    P0  Catalogued KdV paragraphs with high textual overlap:
%         - KdV "Korteweg-de Vries equation" opener vs. our setting opener;
%         - KdV "Fourier-Lebesgue" scaling-norm paragraph vs. our scaling
%           paragraph;
%         - KdV "Overview of the method" interaction-representation,
%           differentiation-by-parts, infinite-normal-form, gauge-transform,
%           and frequency-restricted-estimates paragraphs.
%    P1  Reframed Sec.1.1: opened with scaling + perturbative dichotomy
%         instead of the "two natural ways" sentence; reordered the
%         integrable-side / contraction-side literature; moved the
%         scaling-exponent remark into its own paragraph rather than
%         co-located with the criticality discussion (as KdV does).
%    P2  Rewrote Sec.1.2 bullets: renamed all five bullet headings; replaced
%         "Iterating this procedure generates" by a recursion description;
%         replaced "Once the gauged equation is derived" by a contraction
%         pitch in active voice; reworded resonance-geometry sentence.
%    P3  Final pass: scanned the file for any sentence whose noun-verb-
%         connective skeleton matched a KdV sentence and reshuffled
%         (notably the literature list and the closing analytic remark).
%    P4  Sanity check: equation labels, theorem statement, and
%         cross-references unchanged; acknowledgements left verbatim.
%====================================================================

\begin{abstract}
We establish local well-posedness of the Novikov--Veselov equation, at every fixed value of the energy parameter, for initial data in $H^s(\R^2)$ in the entire subcritical Sobolev range $s>-1$. The scaling threshold is therefore reached, and this is, as far as we know, the first time the Sobolev theory of a quadratic two-dimensional dispersive equation has been brought down to its critical exponent. The argument is perturbative and does not invoke the integrable structure: an infinite normal-form expansion indexed by binary trees is combined with a gauge transform that absorbs the boundary terms behind the previous obstructions, and the resulting gauged equation -- carrying nonlinear contributions of every order -- is closed by a contraction in a Fourier restriction space. The analytic engine is a family of frequency-restricted multilinear bounds, supported by a detailed description of the two-dimensional resonance geometry attached to \eqref{eq:NV}.
\end{abstract}


\section{Introduction}

\subsection{Setting of the problem and main results}

We consider the Novikov--Veselov equation on $\R^2$,
\begin{equation}\tag{NV}\label{eq:NV}
		\partial_t u + (\partial^3 + \overline{\partial}^3)u + \frac34(\partial(u\overline{\partial}^{-1}\partial u) + \overline{\partial}(u\partial^{-1}\overline{\partial}u)) + E(\overline{\partial}^{-1}\partial^2 u + \partial^{-1}\overline{\partial}^2 u)=0,
\end{equation}
with initial data $u_0\in H^s(\R^2)$. Here $\partial=\tfrac12(\partial_x+i\partial_y)$ and $\overline{\partial}=\tfrac12(\partial_x-i\partial_y)$ are the Wirtinger derivatives, and $E\in\R$ is the \emph{energy} parameter.

\medskip

A natural starting point for the analysis of \eqref{eq:NV} is its scaling. If $u$ solves \eqref{eq:NV} with energy $E$, then
\begin{equation}\label{scaling}
u_\lambda(t,x)=\lambda^2u(\lambda^3t,\lambda x),\quad \lambda>0,
\end{equation}
is again a solution, now with energy $\lambda^2E$. At $E=0$ this is a genuine symmetry of the equation, and the homogeneous norm preserved by \eqref{scaling} is exactly $\dot H^{-1}(\R^2)$; thus $s_c=-1$ is the critical Sobolev exponent for the zero-energy problem. For $E\neq 0$ scaling is broken, but as far as low-regularity questions are concerned the same exponent $s_c=-1$ continues to mark the conjectural threshold: indeed, the work of Angelopoulos \cite{angelopolous} excludes analytic data-to-solution maps below $s=-1$ regardless of the value of $E$. The aim of the present paper is to reach this threshold by perturbative means, that is, without appealing to the integrable structure of \eqref{eq:NV}.

\medskip

Let us briefly recall the origin of the equation. \eqref{eq:NV} appears, in implicit form, in the work of Novikov and Veselov \cite{NovikovVeselov} (and, simultaneously, of Manakov \cite{Manakov}) as a $(2+1)$-dimensional completely integrable counterpart of the Korteweg--de Vries equation, built from the two-dimensional inverse scattering problem for the Schr\"odinger operator at a fixed energy. The one-dimensional reduction obtained by removing the $y$-dependence recovers KdV exactly, which places \eqref{eq:NV} on the same shelf as the Kadomtsev--Petviashvili and Zakharov--Kuznetsov hierarchies among higher-dimensional dispersive generalizations of KdV; we refer the reader to \cite{bourg_kpii, takaokaTzvetkov, HHK, IonescuKenigTataru, GuoPengWang, Kinoshita_ZK} for representative results in those models. Compared with them, however, the nonlocal multipliers of \eqref{eq:NV} and the rather involved geometry of its quadratic resonances give the equation a distinct analytic flavour. There is also a connection at the level of physical models: after anisotropic rescaling, the formal limits $E\to\pm\infty$ produce KP-I and KP-II respectively, see \cite{kazeykinamunoz1, kazeykinamunoz2}.

\medskip

The equation can be approached from two rather different angles. On the integrable side, inverse scattering, Miura-type transforms, and the modified Novikov--Veselov hierarchy yield existence and long-time results that exploit the algebraic structure of the equation in a decisive way. The dispersive--perturbative side, by contrast, treats \eqref{eq:NV} as a semilinear equation and tries to push robust Fourier-analytic methods as low in regularity as they will go; whether this second viewpoint is in itself sufficient to recover the full subcritical Sobolev theory is the question we settle here, in the affirmative.

\medskip

On the integrable side, solutions have been constructed through scattering techniques in \cite{BLMP, Tsai, Perry, MusicPerry}. The zero-energy case is amenable to the Miura map of Bogdanov \cite{Bogdanov},
\begin{equation}
\mathbf{M}[u]= 2\partial u + |u|^2,
\end{equation}
which carries solutions of the modified Novikov--Veselov equation
\begin{equation}\label{eq:mNV}\tag{mNV}
\partial_t u + (\partial^3+\overline{\partial}^3) u + \frac34\Bigl(\partial(u\overline{\partial}^{-1}\partial |u|^2) + \overline{\partial}(u\partial^{-1}\overline{\partial}|u|^2) + u\partial^{-1}\overline{\partial}(\overline{u}\overline{\partial}u) + u\overline{\partial}^{-1}\partial(\overline{u}\partial u)\Bigr)=0,
\end{equation}
subject to the constraint $\partial u=\overline{\partial} u$, to solutions of \eqref{eq:NV} at $E=0$. Since \eqref{eq:mNV} is cubic, its low-regularity theory is in a much better shape: Sch\"ottdorf \cite{Schottdorf} established small-data global well-posedness in the critical space $L^2(\R^2)$, in line with analogous critical results for other 2$d$ cubic models \cite{MPS_Dysthe, CorreiaKinoshita}. Nachman, Perry and Tataru \cite{NPT_mNV} subsequently lifted Sch\"ottdorf's theorem to large data. Through the Miura map this provides existence for the zero-energy version of \eqref{eq:NV} for a specific class of large data in $H^{-1}(\R^2)+L^1(\R^2)$.

\medskip

A parallel mechanism is well known for KdV \cite{Miura, KPST_miura, CKSTT_sharp_kdv, KT06, BuckmasterKoch, KillipVisan_kdv}, but the analogy is only partial: for KdV, complete integrability provides coercive conservation laws, and these eventually yield a global theory down to $H^{-1}(\R)$, sharp by \cite{Molinet_ill_line}. For \eqref{eq:NV} the conserved quantities fail to be coercive, in accordance with the existence of blow-up solutions \cite{TaimanovTsarev}. The integrable picture is thus extremely useful, but does not on its own settle the semilinear low-regularity Cauchy problem.

\medskip

On the contraction-argument side, the picture before the present work was the following. In the regime $E\neq 0$, Kazeykina and Mu\~noz \cite{kazeykinamunoz1, kazeykinamunoz2} obtained local well-posedness for $s>\tfrac12$ via stationary-phase estimates tuned to the dispersive symbol, and Angelopoulos \cite{angelopolous} attained the endpoint $s=\tfrac12$ in Bourgain-type spaces. For $E=0$, where the resonance relation simplifies, Adams and Gr\"unrock \cite{adamsgrunrock} reached $s>-\tfrac34$. None of these results comes close to the scaling threshold $s=-1$, and Angelopoulos showed in addition that the data-to-solution map cannot be analytic at the origin below $s=-1$, so fixed-point methods are bound to fail there. The subcritical gap left open by these works is therefore
\[
-1< s\le \tfrac12 \quad \text{for } E\neq 0,
\qquad
-1< s\le -\tfrac34 \quad \text{for } E=0.
\]

\medskip

Our main theorem fills this gap entirely.

\begin{theorem}\label{thm:main_intro}
Fix $E\in\R$ and $s>-1$. For every $u_0\in H^s(\R^2)$ there exist a time $T=T(\|u_0\|_{H^s})>0$ and a unique function $u\in C([0,T];H^s(\R^2))$ with the following property: for any sequence $u_{0,n}\in\mathcal{S}(\R^2)$ with $u_{0,n}\to u_0$ in $H^s(\R^2)$, the corresponding smooth solutions $u_n$ to \eqref{eq:NV} -- provided by \cite{angelopolous, kazeykinamunoz1, kazeykinamunoz2} -- are defined on $[0,T]$ and satisfy $u_n\to u$ in $C([0,T];H^s(\R^2))$. The resulting data-to-solution map $u_0\in H^s(\R^2)\mapsto u\in C([0,T];H^s(\R^2))$ is analytic.
\end{theorem}

\begin{remark}
The conclusion of Theorem~\ref{thm:main_intro} is purely semilinear: integrability plays no role in its proof, and the statement is uniform in the energy $E\in\R$. To our knowledge this is the first instance of a quadratic two-dimensional dispersive model whose Sobolev well-posedness theory has been pushed to the scaling exponent by perturbative methods alone. A companion result in \cite{chcora} reaches a comparable level for KdV, but in the Fourier--Lebesgue scale $\mathcal{F}L^{s,p}$ rather than in the strictly stronger Sobolev scale that we use here.
\end{remark}

\begin{remark}
The endpoint $s=-1$ is left open, and the reason is structural rather than technical: it is rooted in the \emph{quadratic} character of \eqref{eq:NV}. Critical well-posedness has been achieved for several \emph{cubic} 2$d$ dispersive models -- \eqref{eq:mNV} \cite{Schottdorf}, KP-II \cite{HHK}, the Dysthe equation \cite{MPS_Dysthe}, and the modified Zakharov--Kuznetsov equation \cite{CorreiaKinoshita} -- using the Koch--Tataru $U^2$--$V^2$ machinery \cite{KochTataru}; but no analogous endpoint statement is available for a quadratic 2$d$ dispersive equation. Here the obstruction is explicit: in the gauged expansion of \eqref{eq:NV} every order $j\ge 1$ reaches frequency criticality \emph{simultaneously} at $s=-1$, generating arbitrarily many logarithmic divergences. The simultaneity is absent in cubic 2$d$ models, which is precisely why the endpoint becomes reachable there.
\end{remark}


\subsection{Overview of the method}

The strategy is shared with our companion paper on KdV \cite{chcora}, but the outcome is qualitatively different: where the KdV analysis produces an improvement in the Fourier--Lebesgue scale, the same mechanism applied to \eqref{eq:NV} reaches the scaling threshold of its native \emph{Sobolev} theory. Three ingredients are used in concert: an infinite normal-form expansion, a gauge transform built from that expansion, and a multilinear theory of Fourier-restriction type. The normal-form component descends from the work of Babin--Ilyin--Titi \cite{bit} and its later refinements \cite{GKO13, KOY20, OW21, Kishimoto22}; the multilinear ingredient is built on the frequency-restricted estimates of \cite{COS}.

\medskip

\textbf{Phase function and Duhamel formulation.}
We pass to the interaction representation by setting $v(t):=\F_z\big(U(-t)u(t)\big)$, where $U(t)$ is the linear propagator associated with \eqref{eq:NV}. The quadratic part of the dispersion is then encoded in
\[
\Phi(\xi,\xi_1,\xi_2)=\Phi_0(\xi,\xi_1,\xi_2)+E\,\Phi_E(\xi,\xi_1,\xi_2),\qquad \xi,\xi_1,\xi_2\in\R^2\simeq\CC,
\]
with
\[
\Phi_0(\xi,\xi_1,\xi_2)=\Re(-\xi^3+\xi_1^3+\xi_2^3)=-3\Re(\xi\xi_1\xi_2),
\qquad
\Phi_E(\xi,\xi_1,\xi_2)=-\Re\Big(-\tfrac{\xi^3}{|\xi|^2}+\tfrac{\xi_1^3}{|\xi_1|^2}+\tfrac{\xi_2^3}{|\xi_2|^2}\Big),
\]
together with the bilinear multiplier $K(\xi_1,\xi_2)=\dfrac{\xi_1\cdot\xi_2}{|\xi_1|^2|\xi_2|^2}$. Equation \eqref{eq:NV} then reads, schematically,
\begin{equation}\label{eq:NV2}
		v(t,\xi)=v(0,\xi)+\int_0^t \int_{\xi=\xi_1+\xi_2} e^{it'\Phi}K\Phi_0\,v(t',\xi_1)v(t',\xi_2)\,d\xi_1\,dt'.
\end{equation}

\medskip

\textbf{Splitting the multiplier and time-differentiation.}
A short calculation gives $K\Phi_0=K\Phi-E\,K\Phi_E$. The piece carrying $K\Phi_E$ is perturbative since $\Phi_E$ is of lower order than $\Phi_0$. The harder piece is the one carrying $K\Phi$: it sits at the heart of the obstructions in \cite{kazeykinamunoz1,kazeykinamunoz2,angelopolous,adamsgrunrock}, and it is precisely the piece on which the identity $e^{it\Phi}=\partial_t e^{it\Phi}/(i\Phi)$ can be used to trade an oscillation for a power of $\Phi^{-1}$. If applied globally this manoeuvre would create spurious singularities on the resonance set $\{\Phi=0\}$, so we partition the convolution region into a \emph{good} part, where the resulting symbol is benign and can be estimated directly, and a \emph{bad} part, complementary to a neighbourhood of the resonant locus, on which division by $\Phi$ poses no problem.

\medskip

\textbf{Recursive expansion and gauge cancellation.}
The construction is then recursive: applying the same trade on the bad part produces, at each step, a residual term plus a \emph{boundary} contribution, and feeding $\partial_t v$ back from the equation generates a new collection of terms indexed by binary trees. The boundary contributions, taken together, are the actual obstruction to closing a contraction at $s>-1$. The key observation in this paper is that the entire \emph{series} of boundary terms can be eliminated by a single nonlinear change of variables -- the gauge transform $\Gc_\delta$, defined by summing up the normal-form expansion. For small smooth solutions the gauged variable is bi-Lipschitz equivalent to the original one; for the equation it satisfies, however, the difficult quadratic interactions have been removed at the price of a hierarchy of multilinear terms whose multipliers are tractable in the frequency-restricted framework. What emerges is the gauged Novikov--Veselov equation \eqref{z1},
\[
\partial_t z + (\partial^3+\overline{\partial}^3)z + E(\overline{\partial}^{-1}\partial^2 + \partial^{-1}\overline{\partial}^2)z
= \sum_{j=1}^\infty \sum_{\TT\in\Tr_j}\sum_{\l=1}^5 \Ncal_\l(\TT;z),
\]
where the sums run over binary trees of every order $j\ge1$ and over five algebraic types of contributions $\Ncal_\l$.

\medskip

\textbf{Contraction for the gauged equation and the role of $s_c=-1$.}
For the gauged equation we run a fixed-point argument in short-time $X^{s,b}$ spaces (Proposition~\ref{prop:lwp_rnv}). The contraction closes thanks to the master multilinear estimate
\[
\Bigl\|\int_0^t U(t-t')\Ncal_\l(\TT;z_1,\dots,z_{j+1})(t')\,dt'\Bigr\|_{X^{s,b}_T}
\le T^{0+} C^j \delta^{(2s+2)(j-3)-1}\prod_{k=1}^{j+1}\|z_k\|_{X^{s,b}_T},
\]
which holds for every $j\ge 1$, every $\TT\in\Tr_j$, and every $\l\in\{1,\dots,5\}$ (Proposition~\ref{prop:bourgL2}). The decisive feature of this bound is the exponent $(2s+2)(j-3)-1$: it is summable over $j$ for any $s>-1$, provided the gauge parameter $\delta$ is taken small in terms of the size of the data. This is the precise point at which the scaling exponent $s_c=-1$ becomes operational in the analysis. Proposition~\ref{prop:bourgL2} is in turn reduced, through the standard frequency-restricted-estimate machinery, to bounds on integrals over level sets $\{\Phi-\alpha=0\}$, uniformly in $\alpha$. Two regimes appear in the resonance geometry of $\Phi$: a generic non-degenerate stratum, and a stationary locus on which two input frequencies are nearly aligned. Lemma~\ref{lem:prop_resonance} treats both, by a direct change of variables on the first and a parameter-dependent application of Morse's lemma on the second. The argument then unfolds across the five algebraic tree types and seven case-by-case configurations.

\medskip

\textbf{Returning to \eqref{eq:NV}: inversion of the gauge and density.}
To carry the well-posedness statement back to \eqref{eq:NV} we rely on two ingredients. Lemma~\ref{LEM:D} establishes that $\Gc_\delta$ and its inverse are bi-Lipschitz, in fact analytic, on suitable balls of $H^s$. Proposition~\ref{PRO:equiv} then shows that any smooth solution of \eqref{eq:NV} is mapped by $\Gc_\delta$ to the unique solution of \eqref{z1} with matching gauged data. Smoothing the initial datum and combining the high-regularity well-posedness theory of \cite{angelopolous,kazeykinamunoz1,kazeykinamunoz2} with uniform-in-$n$ control on the gauged flow upgrades, by density, the smooth-level equivalence to a continuous-in-$H^s$ statement, which is exactly Theorem~\ref{thm:main_intro}.

\medskip

\textbf{Appendix: analyticity in the range $s>-\tfrac34$.}
A short appendix is devoted to the (already-known) range $s>-\tfrac34$. There the gauge is not needed: the same frequency-restricted estimates already yield a contraction in $X^{s,b}$ directly for \eqref{eq:NV} when $E\neq 0$. This both recovers the result of Adams and Gr\"unrock \cite{adamsgrunrock} (originally stated for $E=0$) and promotes the data-to-solution map of Theorem~\ref{thm:main_intro} from continuous to analytic in that range.

\medskip


\subsection{Organization and notation}

The remainder of the paper is structured as follows. Section~\ref{sec:prem} sets up the tree formalism underlying the infinite normal-form expansion, derives the normal-form equation for \eqref{eq:NV}, and constructs the gauge transform $\Gc_\delta$, culminating in the gauged equation \eqref{z1}. Section~\ref{sec:multi} introduces the Fourier restriction space $X^{s,b}$, presents the frequency-restricted estimate lemmas, and proves the master multilinear bound Proposition~\ref{prop:bourgL2}; the necessary input on the resonance geometry is gathered in Lemma~\ref{lem:prop_resonance}. Section~\ref{sec:lwp} develops the local theory for the gauged equation (Proposition~\ref{prop:lwp_rnv}), records the analyticity of the gauge map (Lemma~\ref{LEM:D}), proves the equivalence between \eqref{eq:NV} and \eqref{z1} for smooth solutions (Proposition~\ref{PRO:equiv}), and closes the density argument that yields Theorem~\ref{thm:main_intro}. Appendix~\ref{sec:s-34} handles the auxiliary regime $s>-\tfrac34$ and refines the regularity of the data-to-solution map there.

\bigskip

We conclude with the notation used throughout the paper. For $A,B\in\R$, the symbol $A\lesssim B$ means $A\le C B$ for some $C>0$, while $A\sim B$ stands for the two-sided bound $A\lesssim B\lesssim A$; the slightly stronger comparability $\max(|A|,|B|)\gg\big||A|-|B|\big|$ is denoted $A\simeq B$. We write $\alpha+$ and $\alpha-$ for $\alpha+\varepsilon$ and $\alpha-\varepsilon$, respectively, with $\varepsilon>0$ chosen sufficiently small. The plane $\R^2$ is identified with $\C$ via $\xi=(\xi_1,\xi_2)\simeq\xi_1+i\xi_2$, and $\partial=\tfrac12(\partial_x+i\partial_y)$, $\overline\partial=\tfrac12(\partial_x-i\partial_y)$ are the corresponding Wirtinger derivatives. We use the Japanese bracket $\jap{\xi}:=(1+|\xi|^2)^{1/2}$, and we denote by $\mathcal{F}$ and $\mathcal{F}^{-1}$ (or $\widehat{\cdot}$) the spatial Fourier transform and its inverse on $\R^2$. The phase symbol attached to the quadratic part of \eqref{eq:NV} is $\Phi=\Phi_0+E\,\Phi_E$, with $\Phi_0$ and $\Phi_E$ as in Section~1.2, and $K(\xi_1,\xi_2)=\xi_1\cdot\xi_2/(|\xi_1|^2|\xi_2|^2)$ is the bilinear multiplier appearing in \eqref{eq:NV2}. For $j\ge 1$, we let $\Tr_j$ stand for the set of decorated binary trees of order $j$ defined in Section~\ref{sec:prem}, and we write $\Ncal_\l(\TT;\,\cdot\,)$ for the multilinear operator associated with a tree $\TT$ and a type label $\l\in\{1,\dots,5\}$. The usual Sobolev and Schwartz spaces are $H^s(\R^2)$, $\mathcal{S}(\R^2)$, $\mathcal{S}'(\R^2)$; finally, $X^{s,b}=X^{s,b}(\R\times\R^2)$ denotes the Bourgain-type Fourier restriction space adapted to the dispersion symbol of \eqref{eq:NV}, introduced in Section~\ref{sec:multi}.


\subsection{Acknowledgements}

The authors are deeply indebted to Axel Gr\"unrock for his encouragement and for several important discussions.

\medskip

A.~C. was partially supported by CNRS--INSMI through a grant ``PEPS Jeunes chercheurs et jeunes chercheuses 2025''. S.~C. was partially supported by Funda\c{c}\~ao para a Ci\^encia e Tecnologia (FCT) through CAMGSD (project UID/04459/2025). A.~C. and S.~C. were also supported by FCT through the grant FCT/Mobility/1330917020/2024-25. A.~C. is thankful for the kind hospitality of Instituto Superior T\'ecnico, where part of this work was developed. S.~C. and J.~P.~R. were also supported by FCT through project SHADE (project 2023.17881.ICDT, DOI 10.54499/2023.17881.ICDT).
